\section{Completion of the forward marker theorem}\label{sec:forward-completion}

The preceding rounds have made the tested classical descendants
irreducible and excluded their direct decompositions.  We now remove
the remaining field and tensor structures.  The order matters: the
first tensor argument excludes factorizations with two large factors
but leaves a few possible retained plane factors.  The field and
exchanged-pair arguments then remove those planes, because each carries
one of the structures tested there.  Once these exclusions are complete,
every nontrivial inner subgroup normalized by a previous-good label
spans the full endomorphism algebra.  This permits the final
tensor-system exclusion and the construction of actual normal simple
descendants.

Work relative to a fixed path ancestor, with almost simple quotient
\(S\le M\le\operatorname{Aut}(S)\), and retain the ordinary subgroup
labels \(K_r\).  Its natural space is \(V\), its ground field is \(F\),
and \(\mathcal A=\operatorname{End}_F(V)\), with \(v=\dim_F V\).
The base already supplements \(S\).  The witness labels \(r_D,g_D\),
the cut labels \(C_D\), and the axis labels \(M_D\) have the relative
meaning fixed in \Cref{fa:relative-labels}.  Every group tested below
contains a previous-good axis label; its omission bound was enlarged
before the present round as in \Cref{fa:all-ancestors}.

The pure-link normal test of \Cref{fa:pure-links} applies with lower
endpoint \(g_D\), or \(g_*\) at a simultaneous intersection.  Missing
normal powers supply the independently supported subgroups of
\Cref{fa:intersection-support}; we do not require them to exhaust the
intersection with that power.  For disjoint witness supports, a proper
increase above their meet must obey the join restriction in
\Cref{fa:disjoint-join}.  These are the interval properties used in the
remaining exclusions.

All lower dimension bounds in this section are fixed absolute
thresholds, chosen before the marker reservoir.  The representation
bounds and coefficient-twist convention are those of
\Cref{fa:degree-gap,fa:coefficient-twists}.  After the structural
exclusions, we select descendants having the additional root and tensor
protections relative to every earlier member of a classical run.  The
absolute path bound will then force the desired point configuration.

\subsection{Overgroups of field and paired stabilizers}

We first establish the overgroup information needed by the field and
paired tests.  Under the following hypotheses, a prime-degree field
stabilizer or an exchanged-pair stabilizer is maximal.  For a field of
two-prime degree, every maximal overgroup must preserve one of its two
prime-degree subfields.  Thus, when two disjoint witnesses have such a
field stabilizer as their intersection, they cannot leave room for the
third cut supplied by \Cref{fa:cut-meet}.

\begin{proposition}[Field and paired overgroups in the prepared configuration]
\label{fb:field-pair-overgroups}
Use the transferred forms of \Cref{sec:classical}.
Let \(D\subseteq\mathcal A\) be a field containing \(F\), with \([D:F]=n\), where \(n\) is prime or a
product of two distinct primes, and let \(m=v/n\) be sufficiently large.
In an isometry type assume that \(D\) is selfadjoint; in orthogonal
characteristic two assume the irreducible normalization already obtained in the
preceding rounds.  Omit symplectic characteristic two with nontrivial adjoint
on \(D\).

Alternatively, in an isometry type other than symplectic characteristic two,
let \(e,e^*=1-e\) be exchanged complementary idempotents, with totally
singular images when a quadratic form is present, and let \(m=v/2\) be
sufficiently large.  Suppose that the collineation action of each overgroup
under consideration is semilinearly irreducible, as ensured here by the
previous-good axis.

Every proper overgroup in \(M\) of the full field stabilizer preserves a
nontrivial subfield algebra of \(D\) over \(F\).  Every proper overgroup of
the full paired stabilizer preserves that pair.  Consequently a prime-degree
field stabilizer and an exchanged-pair stabilizer are maximal; for a
compositum of two prime degrees its maximal overgroups can only be the
stabilizers of its prime-degree subfields.
\end{proposition}
\begin{proof}

In the field case the stabilizer includes the perfect transfer group \(Y\) over \(D\), of large natural degree \(m\). In the paired case use the paired perfect \({\rm SL}_m(F)\). Its centralizer algebra is \(Fe+F(1-e)\), by natural versus adjoint-dual representations; that for the field case is \(D\).

The full stabilizer contains a supplement to \(S\). To prove the assertion take a minimal subgroup \(N\ne1\) normal in the overgroup inside \(S\). The lift algebra is semisimple, and no center components can be permuted nontrivially by \(Y\), by its minimum permutation degree (at least exponential in an absolute positive constant times \(m f_Y\), where \(p^{f_Y}\) is its field parameter with \(m f_Y\ge v/2\); the number of center components is at most \(v\)). Thus its center lies in the displayed centralizer alternatives (perfection). Any enlargement from scalars already suffices: a larger center normalized by the overgroup then gives the indicated field (inside \(D\)) or exchanged pair. Here the high-degree centralizer facts use the transfer/paired representations as in the earlier pruning. Nor can there be a proper nonscalar central simple tensor factor, or multiple simple components of \(N\) tensoring to the full algebra: \(Y\) must individually stabilize such matrix factors in either setting (by order), acting projectively nontrivially on every factor space, as the centralizer algebras are commutative.

Indeed any such nontrivial projective representation of \(Y\) over \(F\) has degree greater than \(\sqrt v\). For the paired group this follows from \(m\); in the field case lift on its universal covering type, with some nontrivial composition factor. The factors are stable as a multiset under twist of field height \(h=\log_p |F|\). Indeed coefficient twisting the lifted representation through this height preserves the projective matrices (taken over \(F\)), hence gives the same representation by perfection on the universal cover, and in particular the same composition factors. If that constituent has orbit length \(d\), its restricted highest weight has at least \(f_Y/(d h)\) nonzero digits, since the twist circulates them modulo field height \(f_Y\) of \(Y\), where \(f_Y/h\ge n/2\). Diagram twisting does not affect the support positions. This costs dimension at least \(d m^{\lceil n/(2d)\rceil}\ge nm/2\), counting the orbit constituents with their Steinberg dimensions (use \(m^{j-1}\ge j\) for \(m\ge2,\ j\ge1\)). Thus two tensor directions are already impossible.

The full-algebra elementary abelian normalizer inside \(S\) is again too small by orders to contain \(Y\). Thus otherwise we have a full-spanning simple \(B=N\), with \(Y\le B\) as before (centralizer inside \(S\) trivial). The order/degree exclusions again force \(B\) high classical of the same characteristic, since \(Y\)'s field-height times natural degree is at least a constant times \(h v\). Let its natural dimension be \(w\ge m\), field height \(f_B\), and put \(t=f_B/\gcd(f_B,h)\). The absolutely irreducible projective representation in \(V\) costs \(v\ge w^t\) by \Cref{fa:coefficient-twists}. Orders compared with \(Y\) give
\[
             h v m\le C f_B w^2
\]
for an absolute constant \(C\). With \(m\) large this forces \(t=1\) (otherwise \(m\le C t w^{2-t}\)), then \(w^2\ge vm/C\). Again the representation must be natural up to twist by the nonnatural lower bound of \Cref{fa:degree-gap}.

We detail series and root checks here. In ambient linear type the field support uses the natural linear type over \(D\) with its embedding-conjugates and is not self-dual. In unitary type we necessarily have odd \(n\) (involution extending the ground involution); the \(n\) embedding twists of the natural unitary type occupy distinct digit positions modulo \(n f\), so again cannot match to form a self-dual module. For a pair in unitary type the natural and conjugate-dual of \({\rm SL}_m(F)\) at distinct digits likewise do not pair to a self-dual module. Signed types match by the previous tests; in the included symplectic characteristic two field case (trivial adjoint on \(D\)), \(Y\) is symplectic over \(D\) and cannot preserve a nonsingular quadratic even with extended coefficients. For example its lift is transitive on nonzero ground vectors, impossible for such a quadratic by evaluation there (also over the binary field, using polarization or a 3-space). The square-form preservation by a perfect group would be exact. Now all supporting constituent nonzero digits lie at multiples of the ambient \(f\) with \(f_Y\) a multiple of \(f\) (also in the indicated paired representations). No proper smaller positive height could fix the action projectively, even up to duality. To see the digit-position convention explicitly, the field-transfer action over the closure uses natural copies over the embeddings of \(D\) extending the ground embedding, hence twists spaced by \(h\), with \(f_Y=hn\) in the linear or trivial-adjoint case and \(hn/2\) in the involutory case. In the latter, either \(h=2f\) in unitary type, or \(h=f\) in bilateral type and the nontrivial involution fixes \(F\), so \(n\) is even there. For a pair one uses \(f_Y=h\) with natural and adjoint-dual (thus conjugate-dual in unitary type). Wrapping may diagram-twist but does not affect the positions modulo \(f\). A fixing ambient root would intertwine the representation with a coefficient twist by its height (possibly with duality) as natural projective matrices, hence on perfect lifts as above, rotating the positions to nonmultiples of \(f\). This rules out a proper root as before and proves the assertion. In particular prime field labels and exchanged-pair labels in these uses are maximal. For a two-prime compositum the only possible maximal overgroups correspond to prime subfields.\qedhere

\end{proof}

\begin{lemma}[Canonical quadratics in symplectic characteristic two]
\label{fb:canonical-quadratic-exclusion}
In the prepared symplectic configuration in characteristic two, replicated
failures whose full stabilizers preserve canonical quadratics with the ambient
polar form are impossible.  This excludes ordered tensor stabilizers with
alternating factor forms, nontrivial-adjoint field stabilizers, and exchanged
pair stabilizers as soon as the common previous-good label is irreducible.
\end{lemma}
\begin{proof}

For two disjoint failures whose stabilizers respect such quadratics (transforming as the bilinear form), irreducibility of a common previous-good label forces the same quadratic by square-difference, yet the two labels join to the top, not preserving it. This applies to tensor stabilizers (factor forms then alternating), to fields with nontrivial adjoint by hermitian trace, and exchanged pairs by their canonical hyperbolic quadratic. One can exclude all such additional quadratic stabilizers as soon as previous irreducibility on the joint support is available, also for use when a later tensor label is known to force one of them (alternation of ordered factors by previous irreducible normalization; for a slot system transitivity will give alternation).

\end{proof}

\subsection{Ordered tensor factors}
\begin{proposition}[Exclusion of two large ordered factors]
\label{fb:both-high-factors}
At the present forward-test stage, a homogeneous family of failures preserving
proper ordered matrix tensor factorizations cannot have both factor degrees
above the fixed large-degree threshold.  In isometry types the factors are
selfadjoint and have the normalized factor forms.
\end{proposition}
\begin{proof}

The factor-form arguments of \Cref{sec:classical} apply. Take a homogeneous pattern on all \(k\)-sets with designated algebra \(A_D\) on the retained side (over \(g_D\); the supported special groups are commuting normals with soluble quotient). For disjoint sets these large supported groups normalize each other; nontrivial intersection would force equality of simple groups, hence algebras and labels. Thus they commute projectively and linearly on perfect matrix lifts, so the designated algebras commute and multiply as independent factors. By coloring one can keep all their degrees small relative to total, e.g. at most \(v^{1/10}\), by using arbitrarily many disjoint sets. Now each such factor stabilizer, and the stabilizer of the product on a disjoint pair, is maximal by \Cref{fa:ordered-factor-maximal} and supplementation of the socle by the base. Thus we have labels \(C_D\) and product label \(C_{D\cup E}\) for disjoint \(D,E\), by the three-coatom argument of \Cref{sec:forward} (the aggregate supported group moves its constituent factors). The aggregate small-side algebra \(A_D A_E\) depends only on the union: among such small sides it is unique for its full stabilizer, by the supporting-group identification of individually preserved tensor factors in the observation. By repartitioning two sets at a time, products for any larger disjoint family likewise depend only on the union.

\noindent\emph{Recovering elementary factors.}

Consequently \(A_D\) tensors into elementary matrix factors \(T_i\) indexed by \(i\in D\), independent of \(D\) (scalars allowed). Here are details for \(k\ge2\). At a fixed \(B\) take external mutually disjoint \(k\)-sets \(B_i\) for \(i\in B\). Repartition using sets \(X_i\) each consisting of \(i\) and all but one entry of \(B_i\), plus the final set of leftovers. Then \(A_{X_i}\subseteq A_B A_{B_i}\); take its coefficient algebra \(L_i(B)\) in \(A_B\) relative to this tensor expression (generate from all coefficients). These coefficient algebras commute with each other since the \(A_{X_i}\) commute and the external factors are independent. They jointly generate \(A_B\), since the leftover-set factor lies entirely in the product of the \(A_{B_i}\)'s (repartitioning the auxiliary sets), and by repartition of the whole union including \(B\), \(A_B\) now lies in the product of the coefficient-generated algebra with the independent external tensor factors. Thus they are central simple tensor factors (commuting generating algebras have scalar centers and no radicals). For another choice with disjoint auxiliaries they commute across choices at different positions by the same reasoning, so coincide at matched positions by centralizers. Compare arbitrary choices through spare disjoint auxiliaries. And for a different \(B\) still containing \(i\), one can use the same \(X_i,B_i\) auxiliary choice, getting exactly the same coefficient algebra (take coefficients into the commutant of \(A_{B_i}\)). This proves the assertion; these factors are selfadjoint when needed (take coefficients also in the adjoint basis, since the algebras used are selfadjoint). They commute across arbitrary positions, hence are simultaneously tensor independent by central simplicity. At most one can be trivial by label distinctness, and each nonscalar \(T_i\) has stabilizer \(M_{\{i\}}\) as usual, since it is \(A_D\cap A_E\) for \(D\cap E=\{i\}\) (two such paired meets of cut labels, with remaining indices disjoint, join to \(M_{\{i\}}\); the stabilizer is proper by the tensor order bound). For \(k=1\) this conclusion is immediate instead.

\noindent\emph{The aggregate supplement.}

Take prescribed large fixed set sizes, on nontrivial positions with degrees again kept small enough relative to total by thinning. For all such sizes product-factor stabilizers on sets \(P\) are maximal (large unequal sides), hence \(C_{B(P)}\), \(\emptyset\ne B(P)\subseteq P\); they are distinct and thus \(B(P)=P\) after homogenizing on a sufficiently large reservoir. Use sizes including \(l,2l\) for sufficiently large even \(l\). For disjoint \(D,E\) of size \(l\), \(P=D\cup E\), write \(J=C_D\wedge C_E, F_0=C_P\). On the aggregate for \(P\) we have the full supported high special projectively, say \(R\), normal in \(K_{F_0}\) and lying inside the actual ancestor simple. There is no loss from possible outer cosets: we claim
\[
                  K_J R=K_{F_0}.
\]
Indeed transport \(R\) as a normal product over initial sites at the product label \(F_0\), and join to \(M_P\), giving the product label of \(K_{M_P}R\) in the exact interval \([M_P,F_0]={\rm Sub}(S_{2l})\). Its meet with the alternating group there is either endpoint on \({\rm Sub}(A_{2l})\), by the earlier exact normal test. If the meet is trivial, the projected normal-multiple label in \({\rm Sub}(S_{2l})\) itself has order at most two, and thus preserves some \(l\)-set there (\(l\) even). This puts \(R\) inside the actual factor stabilizer by the cut labels, impossible on its high aggregate by the tensor order bound. So it contains the alternating group formally, giving the claim by joining with \(J\).

\noindent\emph{Degree equality and the characteristic-two exception.}

In orthogonal characteristic two use instead disjoint unequal large
sizes also chosen as above, with their sum.  Then \([J,F_0]\) is
formally just a cover (Young subgroup with unequal blocks in the
combined symmetric group; any movement of the split yields a crossing
transposition by conjugation).  But \(R\) includes the supported
symplectic group on the alternating aggregate (tensor with a further
alternating space, with canonical total quadratic).  Within \(F_0\)
we have a proper intermediate between the split stabilizer \(J\) and
the upper label by imposing the canonical tensor quadratic on the
aggregate, with properness over the split by its supported high
orthogonal group and tensor orders.  Indeed \(K_J\) stabilizes both
ordered pieces and acts by semilinear similarities on them, respecting
the split's canonical aggregate quadratic (zero on pure split tensors)
with the aggregate polar multiplier.  Not all of the supported
symplectic group respects it.  This excludes the case directly.

In the remaining types use our equal-size index sets.  If the two
actual side degrees differ, \Cref{fa:ordered-factor-maximal} makes
\(K_J\cap R\) maximal and nonnormal in \(R\).  This intersection
is the full ordered stabilizer there: within \(R\), fixing one piece
also fixes its commutant piece on the aggregate.  Since
\(K_JR=K_{F_0}\), intersections with \(R\) determine all ordinary
intermediates, so \([K_J,K_{F_0}]\) would be a cover.  Yet the
formal interval has a proper wreath intermediate.
Apply \Cref{fa:product-cover} at the original sites before the
ancestor quotient.  It would make the lower site projections normal
in the upper ones, and therefore make their intersection with the
actual \(R\) normal as well, contradicting the preceding
nonnormality.  Thus all disjoint \(l\)-products of degrees agree.
Comparing replacements with spare sets makes all elementary degrees
agree; finite coloring aligns their form classes.

\noindent\emph{The final permutation count.}

Now under \(F_0\) the label allowing permutations of the chosen slots in \(P\) has kernel label \(M_P\) for that action. It induces all symmetric permutations. Indeed use the aligned tensor transpositions in the ambient socle, with corrections as follows. For common degree \(d\) the interchange has negative-eigenspace dimension \(\binom d2 v/d^2\) for sign purposes in odd characteristic, which is even when \(d\) is even by the many slots, and also when \(d\) is odd and the total dimension even. Linear/unitary determinant can always be corrected (by minus identity for odd total dimension). In orthogonal odd characteristic with even slots the tensor complement of the two has even dimension and square determinant by additional even slots. The negative-space form determinant modulo squares for the involution is then square by tensoring, giving also spinor membership. For odd slots use isometries on the large complement of all the chosen slots in \(P\) for determinant/spinor corrections: tensoring with identity on an odd-dimensional symmetric factor still gives both independent invariants (diagonalize that factor; determinant unchanged and spinor unchanged up to a determinant-dependent scalar class). The full aggregate group exceeds this permutation stabilizer by orders. Thus above \(M_P\) there are more elements than the entire formal interval to \(F_0\), contradiction.

\end{proof}

\begin{lemma}[The remaining ordered-factor possibilities]
\label{fb:nuisance-planes}
After \Cref{fb:both-high-factors}, the missing side of an ordered-factor
failure has large degree.  Its retained side can only be a linear plane over
\(\mathbf F_2\), an odd-characteristic bilateral orthogonal plane, or a
symplectic plane over \(\mathbf F_2\).
\end{lemma}
\begin{proof}

For remaining ordered-factor failures the missing side must still be of large dimension by the missing-site bound: the normal products of supported groups modulo centers have soluble remaining quotient and at most the two-slot ambiguity, so some missing simple orbit would otherwise have bounded type. (If there are two projective simple components on a side, use a missing transitive orbit thereof; supplementation over \(g_D\) gives the top action, here of bounded pair rank.) Thus it is the retained side that is now small. The component list and arguments of the earlier ordered-factor pruning then apply except for retained nuisance factors: a linear plane over two, an odd-characteristic bilinear orthogonal plane, or a symplectic plane over two. Indeed across a disjoint pattern all other retained components lie mutually in the stabilizer labels, with exactly the same normalization and spanning arguments giving commuting distinct component algebras (one or two per label); the cases giving forbidden direct systems there are already excluded. From arbitrarily many disjoint labels then take many distinct commuting slots. Label distinctness indeed requires many distinct slots since the designated slot set determines the designated algebra. The common intersection preserves each label's set of at most two slots (orbit sizes at most two even on the union), so preserves an ordered aggregate of many with still many unused, impossible by \Cref{fb:both-high-factors} on the combined support. This proves the partial exclusion needed.

\end{proof}

\subsection{Field algebras and paired projections}
\begin{lemma}[Transfer degree at full simultaneous stabilizers]
\label{fb:large-transfer}
For the prime-field witnesses in this round, the transferred natural degree
exceeds any prescribed absolute bound after the earlier missing-site tests
have been enlarged.  The same conclusion holds at the full simultaneous
field or paired stabilizers used below, relative to their lower endpoint
\(g_*\).  Prime-field labels are coatoms.  Two disjoint witnesses whose
fields commute cannot have a field compositum.
\end{lemma}
\begin{proof}

For any prime field algebra label under consideration, the large transfer degree hypothesis of the stabilizer observation holds: the transferred special inside the ancestor simple gives, modulo its soluble center, nonabelian simple orbit product(s) with soluble remaining quotient, one of them necessarily missing by the pure-link normal test, hence of high type. To spell out this use of the standard orders/structures: if the module degree were bounded, total ground dimension large would force the extension and its parameter large; permanently soluble dimensions then give only a soluble stabilizer, and otherwise the natural perfect special modulo center gives the simple(s), with at most the split orthogonal four-space ambiguity. Semilinear field automorphisms, similitudes, extra scalars or anti-actions do not affect quotient solubility. We use the same argument for the transferred groups in simultaneous intersections below \emph{when the intersection is the full simultaneous stabilizer as indicated}, over \(g_*\) there. More explicitly, any center lost on passing from such a normal transferred perfect special group to its simple site(s) is soluble normal and retained; if the simple sites are not transitive take a missing orbit (by the normal test and solubility of the rest above the supported products), with boundedly many sites and the lower pure-link endpoint supplying the orbit action by supplementation. Thus the high-type conclusion of the missing-site argument applies; we may demand the transferred natural degree exceed any needed absolute bound, by the bounded-rank test there. Thus field labels are coatoms by \Cref{fb:field-pair-overgroups,fb:canonical-quadratic-exclusion}. For two disjoint-support instances the intersection \(J\) is covered in each, with also a third coatom above it. In particular if the two fields commute with compositum a field, their intersection is the full normalizer of that field, of two-prime degree (equality would give identical labels); its large transfer and \Cref{fb:field-pair-overgroups} rule out the third coatom. For any such application in the symplectic even case with individual trivial adjoints, the compositum also has trivial adjoint.

\end{proof}

\begin{proposition}[Exclusion of fields, pairs, and nuisance planes]
\label{fb:field-pair-exclusion}
In the prepared forward-test configuration, the prime-field and exchanged-pair
failures are impossible.  The nuisance planes of \Cref{fb:nuisance-planes}
are therefore impossible as well, completing the ordered-factor exclusion.
\end{proposition}
\begin{proof}
\noindent\emph{Linear type.}

In linear type the prime-field proof with retained projective tori from \Cref{sec:classical} thus works: commutation gives a field by the direct-system exclusion; the only possible noncommuting exceptional pair gives a quaternion \(M_2\) factor over three normalized by the intersection, now forbidden by \Cref{fb:nuisance-planes}. We thereby also exclude the remaining linear ordered factors (plane over two gives a canonical quadratic field).

\noindent\emph{Fields with nontrivial adjoint in isometry types.}

In isometry types first take field patterns with nontrivial involution restriction. For two disjoint occurrences use \(R=F[D_i,D_j]\) as in the previous pruning, semisimple with either center field or two exchanged central components. We give the adaptation. If a larger center field gives an invariant central prime \(D'\), then either \(K_i\) normalizes it (forcing \(D'=D_i\) by the perfect transfer and centralizer), or the full simultaneous stabilizer of \(D_i,D'\) equals \(J\) by covering. In the latter, if \(F[D_i,D']\) is a field its large transfer inside \(J\) centralizes \(D_j\), giving commutation. If not, there must again be two exchanged components, from equal prime degrees, and the full intersection includes the paired special linear group over \(D_i\) on the components, normalized there (\(D'\) is likewise scalar over \(D_i\) on each, by the tensor of the two equal-degree fields, a product of copies of \(D_i\)). Its two spaces are dually totally isotropic and totally singular when relevant by previous irreducibility. The quotient over this paired type of group is soluble, so by the pure-link normal test and the missing-site bound its module degree on each component is again large. Indeed on a pair commuting with \(D_i\) use the transferred pairing over \(D_i\), acting by a determinant-one matrix and its inverse adjoint. In the split-algebra use this also respects \(D'\) since on each component the two prime fields act by the same field of scalars. Quadratic restrictions with zero polar on the pair spaces in characteristic two are zero by previous-good irreducibility (else the kernel sum gives an invariant proper nonzero subspace), so pairing preservation suffices there. This special lies inside the ancestor simple by perfection in the nonsoluble large-parameter/dimension range, normalized in the simultaneous stabilizer which acts by paired semilinear similarities allowing exchange (thus soluble remaining quotient; dimension one itself would just give solubility). It too centralizes \(D_j\), and its ground-field centralizer commutes with \(D_i\), as in the earlier paired situation: in the now large degree the natural and adjoint-dual constituents over the embeddings are pairwise distinguished, including across the two sides. The same argument applies directly if the center of \(R\) has two exchanged components, using their pair stabilizer which cannot contain \(K_i\) by the field centralizer; again its full intersection with \(K_i\) must be \(J\).

For center \(F\), \(R\ne\mathcal A\) since otherwise the intersection has soluble image on the fields with soluble kernel, impossible. Thus \(R\) gives an ordered tensor, necessarily with retained nuisance plane side. If \(R\) itself is that plane algebra, it cannot contain two distinct fields of the nontrivial-adjoint kind: on an odd orthogonal plane such a field has the two singular lines as eigenspaces over the closure and is uniquely recovered, and over two the quadratic field in a plane matrix algebra is unique. If its commutant is the nuisance plane instead, \(K_i\) cannot preserve the tensor (large transfer and field centralizer), so the full intersection of \(K_i\) with this ordered tensor stabilizer is \(J\). On the large factor we have the perfect transfer relative to \(D_i\), with again large degree by the missing-site argument (similarity projection on the nuisance side is soluble), supported inside \(J\). Here transfer is on that large factor space for its factor form; in the bilinear characteristic-two setting this is alternating, and the tensor with the nuisance-side identity preserves the canonical total quadratic in orthogonal type. Its centralizer in \(R\) is \(D_i\), but would include \(D_j\), impossible. In this transfer-on-a-factor use, the full simultaneous stabilizer acts on that matrix factor by semilinear similarities respecting \(D_i\), so on its factor space one again has only a soluble remaining similarity quotient over the transferred special. On the nuisance plane factor the whole similarity projection is already soluble. Likewise for the simultaneous paired transfers, after fixing field scalars and the two summands, linear action on one summand and the pairing multiplier determine the action on the other. Thus the soluble-quotient tests indeed use the corresponding supported perfect transfer images inside the simultaneous group, without needing arbitrary individual similitudes there to embed unchanged in the ancestor simple.

Consequently any remaining pair commutes. Two distinct commuting prime fields of the same degree generate a split algebra; for homogeneous adjoint type its primitive component projectors are selfadjoint (relative identifications of the fields commute with their indicated adjoints), contrary to direct-system exclusion. Distinct prime degrees give the field compositum contradiction already proved. This finishes the first prime-field case.

\noindent\emph{Exchanged pairs and trivial-adjoint fields.}

Next exclude exchanged pairs \(e,e^*=1-e\) of idempotents; total singularity when necessary follows as before. Their labels likewise are coatoms by \Cref{fb:field-pair-overgroups}. For the algebra generated by two projections from disjoint instances a larger center field would give a prime \(D'\) commuting with them and normalized by the intersection. It is not normalized by the first label (large paired group centralizer). The full intersection of the first pair stabilizer with \(N_M(D')\) thus equals the pair-label intersection \(J\). Again it includes a paired \(D'\)-linear special group of large degree by the same argument. Centralizing the other projection (perfection), it forces commutation of the two projections by distinct natural and adjoint-dual constituents. But distinct commuting exchanged pairs normalized by the intersection give all four cells (zero cells paired oppositely would instead force the same pair), whose sums on adjoint orbits give a forbidden nondegenerate direct split, just as before. A center with two exchanged components likewise would force that impossibility since its pair commutes with both. For center \(F\), their generated algebra is \(M_2(F)\) by the scalar \((e-f)^2\) argument for two idempotents \(e,f\). This must have nuisance plane type by the ordered-factor exclusion. An odd orthogonal plane gives at most one exchanged pair, and the plane over two now violates the nontrivial-adjoint field exclusion. This finishes paired exclusions, hence nuisance planes too by the canonical pair or nontrivial-adjoint quadratic field there. Remaining trivial-adjoint prime fields now have the same two-label proof as above (center \(F\) immediately impossible since all such ordered factors are excluded).

\end{proof}

\subsection{Permuted tensor systems}
We first record the generation calculation that will replace direct-space
merging when the objects being combined are sets of tensor directions.

\begin{lemma}[Merging compatible tensor directions]
\label{fb:tensor-merging}
Consider compatible sets of tensor directions \(A B\) and \(B C\) in the
prepared classical configuration, with common aggregate degree \(\dim B\)
sufficiently large.  Suppose their tensor-supported natural special groups
are available.  They generate a group containing the tensor-supported special
on \(A B C\).  In orthogonal characteristic two, use the canonical orthogonal
special on every subset with at least two slots, with product polar form and
quadratic zero on fine simple tensors.  Symplectic characteristic two has
already been excluded by \Cref{fb:canonical-quadratic-exclusion}.
\end{lemma}
\begin{proof}

On compatible sets of directions call the exclusive parts \(A,C\) and common part \(B\), also denoting their aggregate spaces by these letters. Assume \(B\) has sufficiently large degree, and the two direction sets have their full tensor-supported special groups. In orthogonal ambient characteristic two we need only the canonical orthogonal specials there and on the desired union when multislotted (still no symplectic characteristic two ancestor here). Then these groups generate so as to contain the union special; nested cases need no argument, so take both exclusive parts nonempty. Work first on \(A\otimes B_0\otimes C\) where \(B_0\) is a hyperbolic summand of \(B\) of defect at most two, or the entire \(B\) in linear type. In \(AB\) take transfers \(1+N-N^*\) with \(N\) directed between signed hyperbolic \(B\)-coordinates \(s\to r\) from distinct pairs, and arbitrary \(A\)-matrix there. These are in the natural special: the adjoint arrow uses the opposite coordinates in exchanged order, so \(X=N-N^*\) has square zero and \(1+X\) is a determinant-one isometry (in odd orthogonal type unipotent of odd order gives trivial spinor norm). For the canonical quadratic in characteristic two the displacement \(X\) has totally singular image in the two \(A\)-fibers at mutually orthogonal isotropic \(B\)-coordinates, by canonical zero on pure tensors; also \(b(x,Xx)=0\) by symmetry and adjunction for its product polar form \(b\), and displacement rank is even. In \(BC\) take similarly \(t\to s\) with arbitrary \(C\)-matrix \(L\) in the transfer arrow. Use three distinct pairs. Their commutator gives the paired transfer \(t\to r\) with product arrow tensoring the two arbitrary matrices, by elementary matrix multiplication; the higher cross terms vanish. Thus we have all such arrows on \(A\otimes C\) by linear spanning. In linear type simply use single unpaired arrows with the same calculation. In individual basis coordinates now, even short paired arrows between nonopposite (distinct-pair) coordinates within one \(B\)-pair region follow by commuting through a spare pair. These standard short transfers give the split special on \(A\otimes B_0\otimes C\) supported there: ordinary elementary/root generation (in odd symplectic type commutators give also the long roots with factor two, and in hyperbolic unitary type the commutators through an intermediate index to the opposite of the initial one give all trace-zero long parameters). In orthogonal characteristic two they give the split type D generators for the indicated quadratic, which suffices. These are assertions on natural matrices (or their projective images), with identity on the complementary summand inside \(ABC\), and acting by tensor identity in all outside directions.

To see the generator step explicitly, the arbitrary \(A\otimes C\)-arrows between distinct \(B\)-pairs are additive under transfer multiplication, and the commutator used to obtain them has just the composite arrow \(N_1N_2\) and its negative adjoint when the three pairs are distinct. Choose individual hyperbolic bases on the paired \(A\otimes C\) fibers over \(B_0\), totally singular on each coordinate side (also for the indicated quadratic). Rank-one arrows between distinct big pairs give the elementary paired updates there, then any two distinct individual pairs can use an intermediate signed coordinate from a spare big pair. Thus all distinct-pair root groups are present. If a long root is needed, commutation through an intermediate index with final index opposite to the initial one yields on that last arrow the composite minus its adjoint (higher products still zero), i.e. \(2c\) in the standard odd symplectic coordinates or \(c-\bar c\) in normalized hermitian coordinates, with arbitrary composite parameter \(c\). This supplies all long parameters even in even-characteristic unitary type. Ordinary root generation of the split special here thus applies (hyperbolic orthogonal uses just type \(D\) arrows, including at characteristic two). In linear type the unpaired calculation gives all elementary matrices analogously.

If a small anisotropic part remains, take an isometric copy of it inside \(B_0\). Using also the perfect special on \(B\) (embedded in a given supported group), conjugate to the special supported instead with \(B_0\) moved onto the orthogonal complement of that copy, by Witt transitivity with corrections (the large orthogonal complement supplies determinant/spinor data as in the subspace arguments). Here hyperbolic splitting of \(B\) uses its bilateral or hermitian pairing; in bilateral characteristic two it is alternating with no remaining part. A large enough hyperbolic \(B_0\) represents an isometric subspace copy of the small anisotropic part in the other form cases (e.g. embed together with its negative in hyperbolic pairs). The supported embedding on \(B\) used for transport lies in a given special by tensoring with identity (form preservation and high-rank perfection). These two resulting large direct spaces within \(ABC\) are compatible with large nondegenerate intersection and full sum, so \Cref{fa:direct-merging} finishes the argument. This proves the tensor merging assertion.

In applying this repeatedly in orthogonal characteristic two, use on every multislotted subset the quadratic with the product polar and zero on the fine simple tensors there. It is also zero for pure tensors across any two-part cut there (apply the canonical construction for that cut and compare on a tensor basis), so column groups or their altered versions supply the required specials. Aligned internal fine permutations carry these choices as required; when transporting whole merged sets of at least two columns the form is canonical on column tensors too, so arbitrary column similarities in those transports cause no problem.

\end{proof}

\begin{proposition}[Exclusion of permuted tensor systems]
\label{fb:permuted-tensors}
After the preceding rounds, neither a full-algebra elementary abelian
normalizer nor a nontrivial permuted tensor-system stabilizer can be a failure
label in the forward test.
\end{proposition}
\begin{proof}
\noindent\emph{Common tensor refinement.}

We now have the same full-spanning property as in \Cref{c:full-spanning} for a nontrivial inner subgroup normalized by a previous-good label. Thus the elementary abelian exclusion again works by scalar commutator pairing and inclusion of all retained elementary groups in the common intersection (normal there). Take then the tensor-system labels allowing permutations, with the very same reductions of exceptional small types as there. In the retained-product case we again have minimal normals in the common intersection recovering the slot systems. Otherwise take two disjoint occurrences, with common intersection \(E\), and missing transitive products with rich top action (thus high slot dimensions by the missing-site bound). The supported normal products in \(E\) supplied by \Cref{fa:intersection-support} contain in their images the same base-core monolith (using the lower endpoint \(g_*\)). Hence the tuple-intersection commutator argument, followed by full spanning, gives a common matrix tensor refinement by intersections assigned to slot pairs. Concretely use \(N_i=\prod_a R_{i,a}\trianglelefteq E\) of singly supported subgroups in the two missing products as supplied above (the projective tensor supports here have one nonabelian simple per slot). In the quotient by the core of \(K_{g_*}\), assign each simple factor of the monolith \(W\) for each \(i\) to some \(\bar R_{i,a}\) containing it, using \([W,\bar R_{i,a}]\) since \(W\le\bar N_i\). Then \(\overline{N_1\cap N_2}\) still contains \(W\) by commutators. Just as before take intersections of the assigned \(R_{i,a_i}\) with \(N_1\cap N_2\), normal there and their images containing the assigned simple by commutators; once again taking their mutual commutators gives its containment in the image of \(R_{1,a_1}\cap R_{2,a_2}\). Thus the tuple intersections generate a nontrivial \(E\)-normalized subgroup and jointly span fully. Distinct tuples have matrix lifts commuting linearly by their distinct slot supports in some system. So their nonscalar lift algebras are commuting central simple tensor factors as before (one per occupied slot pair), recovering each coarse algebra by containment and fullness. Use the same selfadjoint/form observations as before. The fine slots have uniform degree \(h\ge2\) and forms aligned up to scalar by previous-good transitivity; each row or column contains at least two by block primitivity in \(E\) (supplementation supplies transitive top on each system including at least all alternating permutations). Indeed a singleton width would occur throughout that system by transitivity and the other system would then group it into blocks, hence coincide by primitivity and having at least two slots in both, impossible for disjoint supports. In orthogonal characteristic two the fine forms are alternating and the total quadratic is canonical, zero on fine simple tensors; this canonical quadratic is also the one for grouping by \emph{any} bipartition of the fine slots (uniqueness from values on a tensor basis). Here \Cref{fb:canonical-quadratic-exclusion} excludes symplectic characteristic two.

\noindent\emph{A local alteration of columns.}

Use \(u\) the full row-system label intersected with the fine-system stabilizer, a proper decrease (order on any high row slot) above \(E\). Within \(u\) we can interchange two fine slots of a row or cycle three, leaving others in place as slots. Indeed use aligned permutations, commutators if at least three. For a transposition the determinant/spinor arguments given earlier work since there are at least four fine slots, using internal corrections on another slot for odd orthogonal degree (then \(h\) odd, at least three, repeated with odd multiplicity). In characteristic two orthogonal type the tensor swap has even displacement rank by the even \(h\) and at least four slots. Thus one column gets altered by a replacement, remaining in large overlap with its original version (common product degree at least the square root of the column degree).

\noindent\emph{Generation and the join contradiction.}

By \Cref{fb:tensor-merging}, adjoining \(u\) to the column label merges a column with its altered version. If the gained fine direction within another column has large degree, merge immediately with that column; otherwise transport to its other slots by aligned within-column alternating permutations (its fine-slot count then at least three) and merge extensions via their large common column. As with direct systems conjugation under the column label then gives merged pairs on a connected graph of columns by primitivity; overlapping edges (and enlarging unions) merge with at least a large column in common, giving the full special projectively, thus the ancestor simple, contradicting \Cref{fa:disjoint-join}. Thus these last systems too are excluded.

\end{proof}

\subsection{Grounded descendants and uniform thresholds}
\begin{proposition}[Grounding at bounded omissions]
\label{fb:grounding}
After the prescribed rounds, every tested descendant has an actual nested
normal simple supplement with soluble centralizer in the base.  Relative to a
high classical ancestor it is absolutely irreducible and avoids all the
excluded tensor structures.  Relative to an alternating ancestor it is
primitive and transitive unless the desired point configuration has already
been obtained.  Its simple type can be required to exceed any fixed bounded
order or bounded-rank threshold.

The same actual normal simplicity and full spanning apply to a label above a
previous-good axis and below the ancestor when it has the indicated
nonabelian part: either a contained perfect simple group as used here, or its
own pure link satisfying the normal test.  In the first case that perfect
group lies in the resulting normal simple.
\end{proposition}
\begin{proof}

First extract a normal simple inside the ancestor socle, using the
structural conclusions just proved.  In classical type this is the
minimal-normal argument of \Cref{c:classical-conclusion}: elementary
abelian normals have been excluded, while multiple nonabelian simple
factors would give a forbidden tensor system through their commuting
lift algebras.  The remaining actual normal simple spans the whole
matrix algebra.  Its centralizer in the ancestor almost simple group
therefore injects into the soluble outer quotient.  In alternating
type the primitive, non-affine, non-Cartesian and non-diagonal
conclusions give an almost simple group, as in
\Cref{c:alternating-conclusion}.

In either case the tested label is soluble over this simple, also
after pulling back the ancestor's soluble centralizer.  The simple
lifts inside the actual ancestor socle, where the quotient map is
faithful; hence successive simples really are nested subgroups.
For an axis descendant the pure-link normal test has lower endpoint
\(H\).  It puts the soluble centralizer in \(H\) and forces the
simple to be missing: if it were retained, the base-core quotient
would be soluble, contradicting its nonabelian monolith.  Thus the
axis label is supplemented by this actual simple.

For a general label above the previous-good axis, the same extraction
applies when its own pure link or a contained perfect simple group
supplies the nonabelian part.  We need no supplementation assertion
with the original base in the latter case.  Solubility of the quotient
over the extracted simple already forces every such contained
perfect group into it.  This is the conclusion used at the third
coatom in the protection arguments.

Finally, at axis descendants the missing-site product test applies
to their actual simple supplements inside the original \(S\), and
\Cref{fa:bounded-types} makes their types exceed any fixed threshold.
Every link and intersection used here has bounded omissions; the
earlier exclusion bounds were enlarged to make those intersections
previous-good.

\end{proof}

\begin{remark}[Choice of thresholds]\label{fb:uniform-thresholds}

In arranging these rounds the lower high-dimension requirements in the merging, large-factor and projection calculations are fixed thresholds, not functions of the field parameters. For instance the bounded-codimension grouping/collapse of near direct summands needs only an absolute number of transports (separate off bounded small-space dimensions first, and in the high smaller dimension use the full special orbit argument); choose the lower dimension cutoffs there accordingly. When forcing actual missing sites to be large, take the bounded-type threshold in their product tests also large enough for subsequent requirements such as large column dimensions/degrees (and hence large overlaps after moving one cell). At full simultaneous field or paired stabilizer intersections a bounded module dimension would force large field size in their transferred actions by the high total dimension; the resulting missing-site tests still have arbitrarily long pure links by the original residual marker count. Product transfers of actual simples at axis descendants likewise require only that many residual markers, not an a priori classification of all product labels at that simple. The counts of indices/replicas used together depend on fixed tests and the subsequent bounded path/omission requirements, not the sizes of those simple groups. Thus the enlargements backwards and homogeneity relative to path prefixes as at the start suffice here.

\end{remark}

\subsection{Selecting protected extensions}\label{fb:protection-setup}
We now use the actual nested simples and the path-length bound of
\Cref{w:prepared-path,w:path-bound}.  The protections to be imposed on each
continuing high-classical run are a single tensor atom and absence of a
fixing prime root, relative to every earlier member of that run.

Write \(S_i=B_D\) for any earlier member of the run, including the
current parent, and work above a large supply of fresh ordered
positions. For a trial extension at the parent with omission set \(W\supseteq D\) take increments of prescribed bounded sizes, all beyond previous positions. At trial nodes of matching (continuing) characteristic and classical type relative to \(D\), we have a single cycle under \(\sigma\) on tensor atoms as before, by the structure grounding. The orbit, label and counting arguments of \Cref{w:atom-cycle,w:coarsening-count,w:root-count} apply, only now heights up to \(K_{M_D}\) above \(K_{M_U}\) are bounded by
\[
             O(n\log\log(n+3)),\qquad n=|U\setminus D|.
\]

\begin{lemma}[Reduction to singleton protection failures]
\label{fb:protection-count}
With the bounded predicates homogenized relative to the allowed ordered
prefixes, choose trial increments \(a<b\) of matching continuation color and
with \(b\) sufficiently large compared with \((|W|+a)^4\).  If a protection
failure survives at a continuing trial extension, then one obtains arbitrarily
many coatom witnesses on distinct singleton positions, of the same protection
kind, corresponding to distinct primes.  All these witnesses are counted for
one fixed later descendant \(B_U\).
\end{lemma}
\begin{proof}

Use preliminary homogeneity for a bounded number of comparisons: existence of the actual path socles and their types/continuation relative to prefixes, the structure exclusions, and whether a bad protection label for \((D,L)\) exists for a specified ordered sublist support (and of which of the two kinds), suffices. Choose extension sizes \(a<b\) from \(W\) having matching color for continuation from \(W\), with \(b\) very large compared with \((|W|+a)^4\) (use three spaced sizes). If not continuing at \(a\), use that next size with no new run protection. Otherwise take \(W\subset L\subset U\) with these increments, both continuing. For a run ancestor \(D\), any bad prime-coarsening or prime-root label for \(B_L\) (here \(B_L\) denotes the simple at \(M_L\)) occurs among those counted for \(B_U\). Indeed \(B_U\le B_L\) by actual normal simplicity and soluble centralizer as before. Individually invariant tensor slots for \(B_L\) must group atoms for \(B_U\), with just the unique prime congruence coarsening at that prime; and fixing roots still fix \(B_U\). Their stabilizer/normalizer labels for \(B_L\) are proper below \(K_{M_D}\) and contain \(K_{M_L}\). If one such bad case occurs, we can vary the \emph{entire} list \(L\setminus D\) (even the positions from \(W\setminus D\)) within \(U\setminus D\), retaining by homogeneity the same test relative to \(D\) and ordered omitted support. These altered intermediate labels still have the indicated actual data, and their simple groups contain \(B_U\) by the same soluble-quotient argument (the altered \(M_L\) still contains \(M_U\), even when no longer under \(M_W\)). Indeed we are comparing throughout relative to the fixed \(D\), which still comes first in index order. If the nonempty omitted support of the bad label has size at least two this gives at least \((n/(|W|+a))^2\) different supports (vary \(L\setminus D\) over all subsets of that size in \(U\setminus D\); each specified support of size \(z\) can then occur with probability at most \((|L\setminus D|/n)^z\)), violating the combined atom-prime and root count \(O(1+n\log\log(n+3))\) against the stipulated spacing. If singleton it gives arbitrarily many coatom labels omitting distinct single positions \(x\), relative to \(D\), of the same kind, still from coarsenings or roots for this same \(B_U\). A singleton omitted support leaves no choice of side group on it. In the tensor case there is just one possible such full stabilizer per prime in the count for \(B_U\), and in the root case at most two (height and diagram uniqueness), so we still get arbitrarily many distinct primes.

\end{proof}

\begin{remark}[The third coatom]\label{fb:protection-third-coatom}
Take two of the singleton coatom witnesses from \Cref{fb:protection-count}.
The following common setup applies to the two possible kinds.

For two such coatom labels we also have a third distinct coatom over their intersection in the interval under \(M_D\), namely the cut on their position pair. Write \(K\) for this third coatom group in the almost simple ancestor. It lies above \(K_{M_U}\) so has full spanning and structure avoidance there. Thus it has an absolutely irreducible normal simple \(B\le S_i\) containing any perfect subgroup in \(K\cap S_i\) by Schreier and the trivial inner centralizer in \(S_i\). We will use these assertions when \(K\) contains the indicated high perfect group; they follow by \Cref{fb:grounding}. Write \(v\) for the ancestor natural dimension.

\end{remark}

\begin{lemma}[Singleton root failures]
\label{fb:singleton-roots}
The singleton coatom witnesses in \Cref{fb:protection-count} cannot be
normalizers of fixing prime roots.
\end{lemma}
\begin{proof}

In the root case take distinct primes \(r,s\). Their roots fixing \(B_U\) combine to \(\Theta\) with \(\Theta^{rs}=\Phi_i\) and the two given roots as powers as in \Cref{w:root-count}. Now \(Y=S_i(\Theta)=(G_i^\Theta)'\) lies in the intersection. Its presence forces \(B\) high classical in the same characteristic by \Cref{fa:degree-gap}, and the natural dimensions compared both ways force \(B\) to have dimension \(v\) and act naturally up to twist in the ancestor. The root series agrees by the same form checks using the natural \(Y\). In particular if the ancestor here is symplectic in characteristic two, its natural subfield \(Y\) from \(\Theta\) cannot preserve a polar-nondegenerate quadratic even over the closure, excluding orthogonal \(B\) acting naturally there. In natural subfield coordinates its symplectic matrices on the finite cover act transitively on nonzero finite vectors (invariance of a putative quadratic would be exact by perfection); a constant quadratic value on those vectors forces zero polar form (for the binary field use three independent vectors and polar bilinearity). In orthogonal ambient characteristic two the preservation test instead excludes a natural symplectic \(B\) itself the same way. Hence \(B=S_i(\Psi)\) with \(\Phi_i=\Psi^e\), \(e>1\). Applying \Cref{w:natural-root} also to \(\Psi\) fixing \(Y\) gives \(\Psi=\Theta^h,\ eh=rs\). Lifts normalizing \(B\) commute with \(\Psi\) by uniqueness with given outer part. Thus \(K\) already normalizes the simple of one of the two given roots, impossible.

\end{proof}

\subsection{Natural recognition in the tensor-coarsening configuration}\label{fb:tensor-configuration}
The root case supplied a subgroup already acting naturally on the
ancestor module.  For tensor coarsenings the subgroup in the two-cut
intersection instead acts as a tensor product of spaced natural twists.
We must prove that the normal simple group in the third coatom still
acts naturally on that module.  The following argument uses this
particular tensor configuration throughout.

Choose two distinct odd primes \(r,s\), leaving many unused primes at least as large. Put \(d=rs,\ m=v^{1/d}\); thus \(m\) is very large even compared to \(d^4\) by the remaining prime factors in the atom cycle (the \(B_U\)-atom degree is at least two, raised to the cycle length divided by \(d\)). In particular throughout the following bounds we can impose sufficiently large absolute constants and \(m\gg d^4\), by using a sufficiently large fixed number of distinct labels. In their intersection is the high derived fixed group \(Y\) on the \(d\)-slot coarsening, as in \Cref{w:coarsening-count}. On the ancestor module over the closure it acts projectively irreducibly with lift tensoring \(d\) natural modules of dimension \(m\) at twists spaced by the ancestor field height \(f\), allowing duals. Indeed the cyclically embedded fixed group projects on one slot with parameter \(p^{df}\), and the remaining slot actions are its successive coefficient twists by the ancestor map, as described there. Restricted weights have these \(d\) distinct natural (or dual) digits, hence the individual \(d\) slots are exactly its tensor atoms. If the ancestor is symplectic in characteristic two, the two coarsenings already give canonical quadratics (polar the given form, zero on their respective simple tensors). These are coefficient-ground rational by Frobenius invariance, and the respective labels respect them up to the semilinear similarity action: maps preserving a system of matrix tensor factors permute pure tensor lines (apply factor identifications to the induced semilinear matrix-algebra maps; an intertwiner after that is scalar). The quadratics are equal by square difference and absolute irreducibility under \(B_U\), using perfect lifts for exact invariance. Their joining labels cannot then generate the ancestor (the natural symplectic group here does not preserve such a quadratic). We exclude that case here.

Thus, for the remaining argument, \(d=rs\) with distinct odd primes,
\(v=m^d\), and \(m\) is larger than a sufficiently large absolute multiple
of \(d^4\).  The group \(Y\le B\le S_i\) has parameter \(p^{df}\), and
its action on \(V\) has precisely \(d\) natural or dual-natural atoms at
positions \(0,f,\ldots,(d-1)f\).  The notation \(B\) still denotes the
normal simple in the third coatom group \(K\), supplied by
\Cref{fb:grounding}.

Both \(Y\) and \(B\) act irreducibly on \(V\). We first show that
\(B\) is high classical in the same characteristic, and bound the
support of the highest weight for its action on \(V\). Those bounds
then prove irreducibility of a different module: the natural module
of \(B\), restricted to \(Y\). This permits us to lift the inclusion
to algebraic groups and compose rational highest-weight
representations. The ensuing universal cuts and weight-multiplicity
comparisons force the action on \(V\) itself to be natural.

\begin{lemma}[A single sparse restricted weight]
\label{fb:sparse-weight}
In this configuration, \(B\) is high classical in the same characteristic,
with natural dimension \(w\ge m\).  Its action on \(V\) has highest
weight \(p^x\lambda\), where \(\lambda\ne0\) is \(p\)-restricted.
Every active type-A node of \(\lambda\) defines a cut whose smaller side
has fewer than \(2d\) entries.  In a signed type, its dominant word has
fewer than \(2d\) nonzero coordinates and has no spin support.
\end{lemma}
\begin{proof}

Again \(B\) contains \(Y\) and is high classical in the same characteristic (\(v=m^d\) is small compared with cross-characteristic/permutation bounds for \(Y\)). Its natural dimension \(w\) is at least \(m\), by the exact minimum-degree bound in \Cref{fa:degree-gap}. Its own tensor atoms on the ancestor module are again one \(\sigma\)-orbit by tensor-system avoidance. They group those of \(Y\), so if nontrivially split, \(K\) already lies below one of the two given coarsening labels, impossible: a prime congruence coarsening here would have to come from the \(d\)-cycle of \(Y\)-atoms, necessarily at \(r\) or \(s\) and coinciding with the corresponding system used above. Thus its lift there has highest weight \(p^x\lambda\) with \(\lambda\ne0\) \(p\)-restricted. At every active node of \(\lambda\) in A the cut size or its complement has size less than \(2d\) (Weyl orbit with binomial lower bound versus \(m^d\)). In signed types there are fewer than \(2d\) nonzero coordinates of the dominant word, and no spin supports. Indeed with rank \(e\) and \(k\) nonzero coordinates the orbit has size at least \(2^{k-1}\binom e k\), and use \(e\ge(m-1)/2\), \(m\gg d^4\): on \(2d\le k\le e\) log-concavity puts the minimum of the bound at one of the endpoints, both too large. Likewise \(\binom{w}{2d}>m^d\) in A.

\end{proof}

\begin{lemma}[Irreducibility on the natural module]
\label{fb:natural-restriction}
Under the preceding hypotheses, the projective natural \(B\)-module remains
irreducible on the finite group \(Y\) over the algebraic closure.
\end{lemma}
\begin{proof}

Suppose the projective natural module of \(B\) is reducible on the
finite group \(Y\) over the closure.  In linear type this puts \(Y\)
in a proper parabolic.  In a form type an invariant subspace either
gives such a parabolic, gives a proper nondegenerate splitting, or
has a nonsingular radical line in the characteristic-two orthogonal
case.  We treat the parabolic case first, then the nondegenerate
splittings with two positive ranks, and finally the line cases.

A proper parabolic preserves the filtration of every nontrivial
rational highest-weight module defined by a noncentral cocharacter
(on the covering group if needed).  That filtration is nontrivial:
the projective kernel cannot contain the cocharacter torus, so
distinct cocharacter weights occur.  Parabolic containment would
therefore contradict irreducibility of \(L_B(p^x\lambda)\) on \(Y\).
For an invariant degenerate subspace in a form type, its radical
gives singular or isotropic containment; in characteristic-two
orthogonal type take the kernel of the quadratic form on the radical,
leaving only a nonsingular radical line if that kernel is zero.
It remains to exclude the nondegenerate and line cases.  We may
suppress the overall weight scale \(p^x\) in the following comparisons:
both the extremal restrictions and the multiplicities scale by the
same factor, by Steinberg's theorem.

Use the connected part of each indicated stabilizer on points, with covering groups: the projective finite source is perfect; component quotients for nondegenerate splittings here are soluble (independent isometries on complementary summands, each connected up to at most orthogonal component data; one can work on exact form-preserving lifts). Thus a connected covering subgroup above it would also have to act irreducibly. One can take the reduced identity component over the indicated connected stabilizer using the ambient central covering: it surjects onto that stabilizer (finite surjection with finitely many components and connected irreducible target), by a central isogeny, and projective invariance on its action already tests linear subspace invariance. Ordinary nondegenerate splittings use the product of the two connected classical groups on points (up to central quotients/covers, which do not change their Weyl actions), including an orthogonal two-space torus if needed. For a nondegenerate splitting with both ranks positive, maximal tori of the two parts use all ambient coordinates, except possibly one unused in two odd parts. Among restrictions of ambient orbit weights of \(\lambda\), those with zero on any missing coordinate have full maximal Euclidean norm (coordinate projection and highest-weight convexity), yet are not transitive under the two-part Weyl group, since the sparse nonempty support and spare zeros can be differently allocated. For instance one block has room (in torus rank) for the whole support, and then a nonzero coordinate can also be reassigned to the other. This contradicts irreducibility even on the corresponding connected group, whose simple has all maximal-norm weights in one orbit by the same convexity. In these coordinates one takes hyperbolic pairs on each part separately; if both are odd orthogonal (odd characteristic), combine the two leftover directions as one ambient pair over the closure, unused on restriction. The Euclidean norm in the surviving torus coordinates is Weyl-invariant for the product too, and pulling its torus back isogenously does not change this real coordinate calculation.

\noindent\emph{Line stabilizers.}

For a line in odd orthogonal ambient type B (odd characteristic), restriction of tori to D uses the same coordinates, with even-sum root lattice. If \(k\) is the last positive position in \(\lambda\), sparsity gives \(\lambda_k=a\) with \(0<a<p\). Short root lowering gives \(\lambda-\epsilon_k\) as a weight (pairing \(2a\) nonzero in the odd characteristic), already impossible by parity (the difference after scaling is still odd-sum). For a line in type D, including the nonsingular line in characteristic two, use the standard \(B_{e-1}\) stabilizer embedding, dropping the last torus coordinate \(\epsilon_e\). Here is some detail for characteristic two. In a split hyperbolic basis \(z_i,t_i\) take the nonsingular vector \(u=z_e+t_e\). On points the connected special isometry group's line stabilizer fixes \(u\) and injects on \(u^\perp/\langle u\rangle\) into the symplectic group: kernel elements must fix the whole hyperplane by the quadratic, and then only the nonsingular symmetry (odd displacement rank) could remain besides identity, outside the special group. A connected subgroup \(R\) fixing \(u\) has full image there: take the torus on the first \(e-1\) hyperbolic pairs, their orthogonal root groups, and for each \(i<e\) combine the two end-pair root groups of roots \(\epsilon_i-\epsilon_e,\epsilon_i+\epsilon_e\) with identical parameter \(a_0\) (and similarly in the opposite orientation). Their product sends \(z_e,t_e\) to \(z_e+a_0 z_i,t_e+a_0 z_i\) respectively and \(t_i\) to \(t_i+a_0u+a_0^2 z_i\), fixing the rest. This gives the long symplectic root groups on the quotient by squaring (surjective on points); all other required roots come from the orthogonal ones already used. Take the connected subgroup generated (with Zariski closure). Thus \(R\) has the whole stabilizer's points and dimension of that symplectic group, and is semisimple (a connected solvable normal would have trivial image). Relative to the indicated maximal torus the generators just given have roots \(\pm\epsilon_i,\ \pm\epsilon_i\pm\epsilon_j\) on \(e-1\) coordinates (nonzero differentials also for the combined roots), exhausting the dimension. The displayed torus is maximal by the finite kernel and rank on the symplectic quotient. So this is type \(B_{e-1}\) with the claimed coordinate restriction. We can use its simply connected cover lifting into the ambient cover (central isogeny lifting as in the path proof). In odd characteristic this is the ordinary odd-dimensional special orthogonal stabilizer on the complement up to components/centers.

For the line comparisons these character coordinates on restrictions can be used on the simply connected covers: they come from the vector weights on the indicated surviving pairs. In particular for the special stabilizer just described the characters \(\epsilon_i,\ i<e\), on its torus pull back as the short-root coordinates of type \(B_{e-1}\), and \(\epsilon_e\) becomes zero even on the cover (its vector pair is fixed by this torus). Lifting from a simply connected covering source through an ambient central covering agrees projectively and hence for these linear vector matrices by absence of scalar characters.

Again let \(k\) be the last positive position and put
\(a=\lambda_k\), so \(k\ll e\) and \(0<a<p\).  Root lowering
with pairing \(a\) gives two distinct ambient weights.  On the
line stabilizer they have the same restriction, since
\(\epsilon_e\) restricts to zero:
\[
 \begin{aligned}
 p^x(\lambda-\epsilon_k+\epsilon_e)
   &\longmapsto p^x(\lambda'-\epsilon_k),\\
 p^x(\lambda-\epsilon_k-\epsilon_e)
   &\longmapsto p^x(\lambda'-\epsilon_k).
 \end{aligned}
\]
Here \(\lambda'\) is the restricted highest weight; extend a
positive-chamber cocharacter by zero in the last coordinate to
verify that it is highest.  The two ambient weight spaces are
independent, so this restricted weight has multiplicity at least
two.  We show that an irreducible module of highest weight
\(p^x\lambda'\) has multiplicity at most one there.

After removing the scale, the relevant top Levi-face module has
highest weight \(a\omega_1\) in tail type B, since the difference
down by \(\epsilon_k\) uses just those tail nodes.  The simple
\(L(a\omega_1)\) is a composition factor of the degree-\(a\)
symmetric power of the natural B Weyl module, also in characteristic
two.  The ordinary Weyl character of that natural module has precisely
the signed coordinate pairs and one zero, independently of
characteristic.  Its symmetric power has a unique top weight, giving
the stated constituent, and multiplicity exactly one at
\((a-1)\epsilon_1\): one zero-weight term and \(a-1\) copies of the
first coordinate are necessary.  This bounds the multiplicity in
the simple constituent by one and proves the contradiction.

\end{proof}

\begin{lemma}[Rational lifting and the composed digits]
\label{fb:rational-lift}
The restriction in \Cref{fb:natural-restriction} is the restriction of a
parameter-restricted rational irreducible of the simply connected algebraic
group of \(Y\), of highest weight \(\eta\).  It gives a morphism into the
simply connected algebraic group of \(B\), preserving the required forms.
The rational composition with \(L_B(p^x\lambda)\) is irreducible.  After
removing the common scale \(p^x\), its highest weight \(\mu\) has a
single natural or dual-natural coefficient one at each nonzero digit.
\end{lemma}
\begin{proof}

This natural restriction lifts as a restricted irreducible for the parameter
of \(Y\), with rational highest weight \(\eta\). It gives a morphism
from the simply connected algebraic group of \(Y\) into that of \(B\),
as in \Cref{w:rational-composition}. The invariant bilateral form,
when relevant, is rationally invariant by self-duality and uniqueness.
For the quadratic in characteristic two, if \(\eta\) has multiple
digits, their self-dual factors have alternating forms and hence supply
the canonical quadratic, equal to the given one by uniqueness.
For a single digit, untwist to a \(p\)-restricted weight using the
standard pinned choices; split Frobenius permutes the finite points.
The cocycle for the finite-invariant quadratic then takes values in
the square module of the dual, a nontrivial \(P_Y\)-restricted simple
module because \(P_Y=p^{df}>2\). By \Cref{w:h1-injective} its rational
cohomology class is zero. Thus the cocycle is a coboundary, and that
coboundary is zero because this module has no finite-fixed vectors.
This proves rational invariance.

The composed rational module is irreducible by finite irreducibility: on the finite covering points the lift into the ambient simply connected group of \(B\) differs at most centrally from lifts of the original finite projective inclusion (same projective natural matrices by construction). Thus on \(L_B(p^x\lambda)\) one still obtains projectively the embedding action just specified with its \(d\) atoms. Its highest weight, after removing the scale \(p^x\), say \(\mu\), has only single natural (or dual-natural) coefficients one at its nonzero digits. Indeed every rational digit factor restricted after twist to the finite group is one digit there and by individual tensor invariance must be an aggregate of the \(d\) twisted natural atoms, thus exactly one by restricted-weight uniqueness (an aggregate of two or more has natural/dual digits at that many distinct positions of the finite restricted tuple). Removing the scale is integral by torus restriction.
\end{proof}

\begin{lemma}[Small cuts force endpoint support]
\label{fb:endpoint-cuts}
In this same configuration, \(\lambda\) is supported only at the two
endpoints in type A, or at the first node in a signed type.
\end{lemma}
\begin{proof}

Here is a cut argument to finish. Restrict natural torus weights of \(B\) to get the multiset \(\Omega\) of the module with highest \(\eta\). As in \Cref{w:universal-cuts}, one fixed Weyl translate of \(\lambda\) maximizes throughout the positive chamber for \(Y\), simultaneously on each active fundamental summand. The restriction of each such translated fundamental weight is itself dominant (its orbit maximum, also against \(Y\)-Weyl translates). Indeed Weyl actions on the image torus extend on it via ambient Weyl elements: after transport of the ambient maximal torus by an image normalizer element, conjugate it back in the centralizer of the image torus. Recall the simultaneous maximum exists since at an interior \(Y\)-cocharacter a maximizing orbit weight of the upper module must actually restrict to its unique highest on composition (weight convexity), hence maximizes on the whole positive chamber. Each active fundamental summand is then maximal there since for each cocharacter there is a common ambient Weyl element maximizing all summands. Since their positive integral sum is \(\mu\), each has support only at endpoints for A-type \(Y\), or the first node for signed \(Y\). Each corresponds to a strict upper cut in \(\Omega\) as in that result: all selected elements dominate all unselected and there are no identical vectors across the boundary, by highest multiplicity one (in signed type all active nodes here are early with spare zero coordinates). The sum of the selected weights is that fundamental restriction. Take the small side of size less than \(2d\); for a co-small A-cut this means the negatives of the unselected values, still with that same sum and strict dominance in the dual multiset instead. More explicitly an exchange across a strict coefficient boundary compares natural weights against each other throughout the \(Y\)-chamber. Equal-vector ties across it would produce two distinct ambient orbit weights mapping to the highest vector weight, contrary to multiplicity one. In signed type one similarly flips nonzero coefficient coordinates in checking the maximum (compensating with zero-coefficient coordinates in D), so the positive directed choices lie nonzero in the positive cone and one can make the same strict comparisons in the full natural multiset.

The middle-coordinate Weyl group in type A for \(Y\) (permuting positions other than first and last), or signed Weyl group away from the first coordinate, fixes this sum, thus preserves the selected multiset by strict domination across the cut. By its small size, every selected weight must \emph{individually} be fixed (nontrivial orbits cost at least \(m/3\)). (A permuted selection with the same sum still attains the maximal \(Y\)-chamber interior pairing, hence is the same multiset.) This includes the highest \(\eta\) or its dual, thus supported also only at the indicated endpoints or first node as appropriate. Its active simple reflections are then outside the cut, not fixed individually there. Dominating those forces each selected weight's difference down from highest to be supported only at the active node of any such reflection, which by individual fixedness allows no difference at all: the corresponding simple root (endpoint, or first in signed type) is not fixed by the subgroup either. Thus the cut takes just the highest line. We conclude \(\lambda\) has only endpoint coefficients in A, or just the first coefficient in signed types.

\end{proof}

\begin{lemma}[Multiplicity forces the natural representation]
\label{fb:natural-recognition}
In the tensor-coarsening configuration above, \(\lambda\) is a natural
weight, up to duality.  Consequently \(w=v\), and \(B\) acts naturally
on \(V\), up to Frobenius twist.
\end{lemma}
\begin{proof}

It cannot have both endpoint coefficients \(a,b>0\) (ambient for \(B\) then A), giving \(\mu=a\eta+b\eta^*\). At the lowest nonzero digit exponent \(y\) of \(\eta\), if \(Y\) is type A with just one active endpoint at that digit, \(\mu\) has both active there, impossible. In fact the two coefficients at that position are nonzero modulo \(p\) by \(0<a,b<p\) and the nonzero lone endpoint digit, with no carry into that position from below. Otherwise choose torus coordinates at the full maximum with first two values \(\eta,\eta-p^y\alpha_1\), last two values \(-\eta^*, -\eta^*+p^y\alpha_1\) (last highest-coordinate minimum first in this listing). These distinct weights occur by digit activity (also for signed type with \(\eta^*=\eta\)), and middle coefficients of \(\lambda\) are constant. Here \(\eta+\eta^*\) is not proportional to \(\alpha_1\) by the support already found in high rank, so the four entries are available distinctly at the maximum, rearranging the middle coordinates. Lowering at each of the two endpoint nodes of \(\lambda\) gives two distinct weights colliding at \(\mu-p^y\alpha_1\), up to the overall scale, impossible by simple-root multiplicity one bounds in the simple composed module.

Thus up to duality \(\lambda=z\omega_1,\ \mu=z\eta\). If \(1<z<p\), a coefficient one at any active node of the lowest \(\eta\)-digit would already contradict the required digit shape of \(\mu\). For coefficient at least two there we have three distinct nonzero nonopposite weights \(\eta,\eta-p^y\alpha,\eta-2p^y\alpha\) for that simple root (again the endpoint/first-node shape in high rank ensures nonproportionality to the root). Assign these to the first three natural coordinates of \(B\) at the maximum (in signed types choose signs using also spare zero-coefficient coordinates). The upper \(A_2\)-face for \(\lambda\) gives the two symmetric-power weights lowering twice (one step to the third coordinate, or twice using the second) colliding again along just \(\alpha\). This forces coefficient one for \(\lambda\).

All the collisions just used are comparisons in the rational composed
simple module.  With scales displayed, the first pair of distinct upper
weights restricts to the same weight
\[
             p^x\mu-p^{x+y}\alpha_1,
\]
and the symmetric-power pair in the last case restricts to
\[
             p^x\mu-2p^{x+y}\alpha.
\]
Each lies below the highest weight in a single simple-root direction, where
the divided-power PBW bound gives multiplicity at most one, as in
\Cref{w:universal-cuts}.  Hence the only remaining coefficient is one.
\end{proof}

\begin{proposition}[Singleton tensor failures]
\label{fb:singleton-tensors}
The singleton coatom witnesses in \Cref{fb:protection-count} cannot be
prime tensor-coarsening stabilizers.
\end{proposition}
\begin{proof}

By \Cref{fb:natural-recognition}, \(w=v\) and \(B\) acts naturally
up to twist. Its series matches the ancestor by the previous tests:
in type A the finite \(Y\) digit tuple is non-self-dual, and the
symplectic characteristic-two ancestor case has been removed.
The third coatom is proper and contains the supplementing base, so
\(B<S_i\). By \Cref{w:natural-root}, a proper root of smaller
height fixes \(B\), hence \(Y\).

Fixedness identifies the projective \(Y\)-action with its coefficient
twist through a height \(0<g<f\), possibly with duality, using the
matrix realization of the ambient point maps. These perfect lifts
agree linearly. But the support of the finite restricted tuple is
\[
 \{0,f,\ldots,(d-1)f\}\subseteq\mathbf Z/(df)\mathbf Z.
\]
A shift through \(0<g<f\) moves every one of these positions to a
nonmultiple of \(f\).  A diagram action on wrapping may exchange a natural
weight with its dual, but cannot change which digit positions are nonzero.
Thus no such smaller-height root fixes \(Y\), giving the required
contradiction.
\end{proof}

\subsection{Conclusion of the forward test}
\begin{proposition}[Completion of the forward marker argument]
\label{fb:forward-completion}
The forward marker assertion of \Cref{r:forward-marker} holds.  In the
prepared configuration, the assumption of no point success at all the
prescribed bounded prefixes is impossible once the initial marker set is
sufficiently large.
\end{proposition}
\begin{proof}

Fix the requested number \(m\) of point labels.  Choose the high-type
thresholds first, and let \(D_0\) be the resulting absolute bound in
\Cref{w:path-bound}.  Reserve more than \(D_0\) path nodes.  Choose
the trial omission sizes successively, with the spacing required by
\Cref{fb:protection-count}; each is a finite function of the bounds
already chosen.  This fixes a bound on every path-prefix size.
The earlier exclusion rounds can now have their further-omission
bounds enlarged backwards to cover every bounded union used by a
later round.  All such comparisons are relative to one fixed
ancestor, as in \Cref{fa:all-ancestors}.

Finally choose the initial marker reservoir large enough for those
finite Ramsey requirements, for \(m\) point labels whenever
\Cref{fa:subset-points} succeeds, and for all residual pure links and
missing-site tests.  A terminal spare reservoir, and spaced prefixes
when needed, provide the fresh indices throughout.  Only the bounded
case and comparison predicates already specified are colored; these
choices do not depend on group dimensions or field sizes.

Start with the actual simple \(S\) in the ordinary site interval.
At each subsequent bounded prefix the construction either supplies
the required \(m\) point stabilizers, or continues with an actual
nested simple supplement by \Cref{fb:grounding}.  In the second
case the socles satisfy the transitivity and full-spanning properties
relative to every ancestor, and \Cref{fa:bounded-types} excludes all
bounded types.  A continuing classical run receives both protections
by \Cref{fb:protection-count,fb:singleton-roots,fb:singleton-tensors};
an alternating or otherwise noncontinuing step needs no new run
protection.  Thus absence of point success at the bounded prefixes
would give more than \(D_0\) actual simple nodes satisfying all the
group-theoretic hypotheses of \Cref{w:path-bound}, a contradiction.
The construction must stop with point success at one of those
prefixes.  Its prefix bound and reservoir threshold depend only on
\(m\) and the fixed choices, proving the asserted uniformity.

\end{proof}
